\section{Selective automation must be valued over requested decisions}
\label{exec14:l8-development}

An assisted workflow chooses both an answer and whether to act on it. These
choices change the population on which accuracy is reported. Let $R_i=1$ mean
that request $i$ receives an automated recommendation and $R_i=0$ mean that it
is referred to a fallback. Let $a_i$ and $b_i$ denote the respective actions,
$\theta_i$ the decision-relevant state, and $\ell(a,\theta)$ the loss relative
to an explicitly specified feasible comparator. If $C_i$ includes all work
before referral, $H_i$ the additional fallback cost and $D_i$ the loss caused
by delay, a complete per-request accounting variable is
\[
 Z_i=C_i+R_i\ell(a_i,\theta_i)
       +(1-R_i)\{H_i+D_i+\ell(b_i,\theta_i)\}.
\]
Loss is monetary here only because the example stipulates a monetary mapping.
Legal feasibility, access to assistance and other constraints can instead
restrict the allowed actions; they should not be silently converted into an
invented dollar penalty. A comparison estimates the change in $Z_i$ over the
same requested population, including requests on which either workflow fails.
An unanswered request has a subsequent action or an unresolved consequence,
not a disappearing row. \citep{exec14_l8_economics}

Consider 100 generated offer decisions: 80 with a USD~100 loss for a wrong
action and 20 with a USD~10,000 loss. Workflow A answers all 80 low-consequence
requests and ten high-consequence requests. It makes one error in each group.
Workflow B answers 60 low-consequence and ten high-consequence requests, making
three errors in the first group and none in the second. A's conditional answer
accuracy is $88/90$, higher than B's $67/70$. These are invented counts that
define the example, not estimates of either a model or a bank's error rate.

Assume A costs USD~2 and B costs USD~4 for every requested decision. Each
referral costs another USD~40 and resolves the action correctly in this
particular construction. Delay costs USD~5 for a low-consequence referral and
USD~50 for a high-consequence referral. The full accounting is
\[
 \begin{aligned}
 Z_A^{\mathrm{total}}&=200+10{,}100+400+500=11{,}200,\\
 Z_B^{\mathrm{total}}&=400+300+1{,}200+600=2{,}500.
 \end{aligned}
\]
B is less accurate among its answers and automates fewer requests, yet has
USD~8,700 less total loss and cost in the constructed portfolio. The fallback
is assumed correct, so these totals are optimistic about review; allowing
review errors adds their losses through the same equation. Equally, the
construction does not imply that more referral is always better. Raising its
cost, delay or error rate can reverse the result. The point is that automation
coverage and conditional accuracy do not determine economic value even when
both have been measured correctly. \citep{exec14_l8_economics}

The review decision can be made explicit. Suppose an unreviewed recommendation
has error probability $q$, a reviewed recommendation has probability $r$, an
error loses $L$, and review plus delay costs $c$. Under this binary-loss model,
review improves expected value precisely when
\[
 (q-r)L>c.
\]
For stipulated $q=.01$, $r=.002$ and $c=20$, the loss threshold is USD~2,500.
Reviewing a USD~100-consequence decision has net value minus USD~19.20;
reviewing a USD~10,000-consequence decision has net value USD~60. The same
apparent confidence therefore supports different review choices. A language
model's stated confidence does not supply $q$, and agreement between model
and reviewer does not identify $r$. These probabilities concern correctness
conditional on the target population and routing rule; they require suitable
labels and a separately reviewed evaluation. If review can introduce errors,
$r$ may exceed $q$. \citep{exec14_l8_economics}

CausalChat supplies a reason to study the presentation and review components
together. Its six-student study reports mean edge-editing times of 3.3 minutes
with LLM text and 0.9 minutes with CausalChat Lite. The study used a fixed
sequence of stages and participants chose which edges to resolve. The paper
also reports concerns about missing citations and the reliability of scientific
support. Those observations motivate a banking comparison of review burden,
target errors and evidence inspection. They do not identify a transferable
productivity effect or justify multiplying 2.4 minutes by a bank salary and
calling the result a saving. A prospective comparison should control task
difficulty, learning and reviewer allocation, while preserving the time spent
on consequential mistakes and unresolved cases. \citep{v12_la_2410_14146}

\section{Retry and fallback costs depend on the failure mechanism}
\label{exec14:l8-retry}

Agent workflows often retry after malformed output, unavailable tools or a
failed scientific check. Let $T$ be the first attempt that passes a specified
check and let $m$ be a fixed attempt cap. The number of attempted calls is
$N=\min(T,m)$, with $T>m$ when none passes. For integer $N$,
$N=\sum_{j=1}^{m}\mathbf1\{N\geq j\}$, giving
\[
 \mathbb E N=\sum_{j=1}^{m}\Pr(T\geq j).
\]
This identity needs no independence assumption. With an identical per-attempt
cost $c$ and terminal fallback charge $F$, expected execution cost is
$c\mathbb E N+F\Pr(T>m)$. Attempt-specific review and transfer costs can be
included inside the corresponding indicators. Passing an operational check
is still separate from answering the causal question correctly; a protocol
that detects only malformed JSON needs an additional scientific-error term.
\citep{exec14_l8_economics}

Two stipulated mechanisms have the same 20\% first-attempt failure rate.
In the first, failures are independent with probability $1/5$. At $m=3$,
\[
 \mathbb E N=1+\tfrac15+\tfrac1{25}=\tfrac{31}{25},
 \qquad \Pr(T>3)=\tfrac1{125}.
\]
In the second, one tenth of requests have a permanent defect. Other requests
have independent transient failure probability $1/9$ per attempt. For $k\geq1$,
the probability that the first $k$ attempts all fail is
\[
 \tfrac1{10}+\tfrac9{10}\left(\tfrac19\right)^k.
\]
The first-attempt failure rate is again $1/5$, but
$\mathbb E N=1+1/5+1/9=59/45$ and the terminal failure probability is
$41/405$, approximately 10.12\%. At USD~2 per attempt and USD~50 for the
terminal fallback, expected costs are USD~2.88 and approximately USD~7.68.
Full enumeration of permanent status and all three transient outcomes gives
the same values as the tail-sum calculation. \citep{exec14_l8_economics}

A missing identification method, inconsistent units, inaccessible evidence or
an ambiguous causal target can behave like the permanent component. A new
prompt does not necessarily change the missing information. Independent
retry arithmetic would substantially understate terminal workload in that
case. The proposed agent should classify the failure, retain the failed
request and route to an appropriate remedy: repair a format, select a supported
method, seek evidence or return an unresolved decision. A retry budget limits
spending; it does not establish that repeated reasoning supplies independent
evidence. The classification itself should be checked because a mistaken
``transient'' label can generate costly repetition.

CausalDS reports marked differences in tool-use strategy and total tokens,
alongside tool-output-management failures and distinct content and abstention
outcomes. Its reported ATE interval coverage uses only seven eligible cases
per model, or five for one model, in that particular exam. These scoped results
support retaining trajectories and component outcomes. They do not supply the
permanent-failure fraction of a banking workflow, and total tokens do not
determine dollars without dated input/output/cache prices and the relevant
billing rules. The prototype should measure the actual sequence of attempts,
which input changed between them, and the eventual fallback. No named service
price or bank runtime is assumed here. \citep{v12_la_2607_08093}

\section{Causal ambiguity needs a declared decision rule}
\label{exec14:l8-ambiguity}

The funding example in Section~\ref{exec13:funding-journey} maps a supplied
response set through a stated profit function and obtains net values from
minus USD~625 to plus USD~1,625. Holding the rate fixed has zero incremental
value in that example. To examine the decision rule, temporarily treat those
endpoints as a declared ambiguity set containing the true state; this is a
conditional decision construction. If the set came from a confidence procedure,
its coverage uncertainty remains an additional issue. It is not a posterior
probability law over the possible profits.

A robust net-value rule holds because the worst net value from acting is
negative. A deterministic minimax-regret rule acts. To see the difference,
define regret in state $\theta$ as
\[
 \mathcal R(a,\theta)=\max_{b\in\{\mathrm{hold},\mathrm{act}\}}
          V(b,\theta)-V(a,\theta).
\]
Acting has worst regret USD~625, whereas holding has worst regret USD~1,625.
Robust value protects the attained payoff; minimax regret protects the shortfall
relative to what could have been chosen with state knowledge. Neither rule is
deduced from the graph or interval alone. They express different preferences
for the same scientifically unresolved funding question.

If randomizing the entire action were an admissible policy, choose action with
probability $p$. Expected regret within the adverse and favorable states is
$625p$ and $1625(1-p)$. Their maximum is minimized at
\[
 p=\frac{1625}{625+1625}=\frac{13}{18},
 \qquad \mathcal R_{\max}=\frac{8125}{18}\approx451.39.
\]
This is a randomized decision across instances, not permission to replace the
rate change by $13/18$ of its size. The latter needs a response model for the
new intervention. A bank may prohibit randomized treasury actions, so the
deterministic analysis can remain the relevant one. An assistant should explain
which criterion generated a recommendation and preserve the alternatives.
It should never silently convert an interval midpoint into a uniquely justified
choice. \citep{exec14_l8_economics,exec13_funding_exact}

Sequential servicing adds a second information requirement: the joint values
of the feasible policies. Suppose two generated worlds give net policy values
$(V_A,V_B)=(-100,-200)$ and $(100,0)$. Policy A exceeds B by USD~100 in both
worlds. Separate coordinate projections give $V_A\in[-100,100]$ and
$V_B\in[-200,0]$, whose unrestricted difference for B minus A spans $[-300,100]$. That
box suggests a possible advantage for B by combining coordinates from
incompatible worlds. It is a conservative enlargement, but it can conceal a
strict policy ranking. A useful policy interface should preserve the joint
world or region identity where available and state when only coordinate
bounds have been computed. An LLM can explain this distinction; it cannot
recover lost joint information from two displayed intervals.
\citep{exec14_l8_economics}

For discovery, the information purchase should concern an unresolved feature
that can change an action, such as a consequential orientation or an admissible
intervention regime. Improving unrelated edges may have no value for the
specified offer. For observation construction, the consequential distinction
may instead be the origin map, inclusion design or label law. A timestamp
audit can be useful while leaving those questions untouched. The four journeys
therefore share an accounting language but require different information.
Their proposed acceptance cases keep unsupported estimators, equal-law
counterexamples and unresolved scientific premises visible.
\citep{exec14_l8_design}

\section{A public prototype buys an option to make a later investment}
\label{exec14:l8-option}

A proposal need not know bank ROI before it can justify a small prototype.
It does need an account of which later decision the prototype could change.
Let $G$ be the eventual net value of constructing a selected capability, after
all subsequent implementation and operating costs in a declared horizon.
The choice today is to build or decline. A pilot supplies a signal $S$, costs
$K$ today and delays the later value by a discount factor $d\in(0,1]$.
Under an explicitly stipulated joint probability model, the two option values
are
\[
 W_0=\max\{0,\mathbb E G\},\qquad
 W_{\mathrm{pilot}}=-K+d\,\mathbb E[\max\{0,\mathbb E(G\mid S)\}].
\]
The maximum inside the expectation represents the option to decline after
an unfavorable signal. It must correspond to a feasible stop decision: costs
already committed before the signal cannot be removed by a retrospective
choice. Pilot cost is excluded from $G$ because it is subtracted separately.
The law of $S$ describes how the prototype informs the commercial and scientific
state; a benchmark score is not automatically such a signal.
\citep{exec14_l8_economics}

For an exact planning example, let $G$ be USD~100,000 or minus USD~60,000 with
equal probability. Immediate construction has expected net value USD~20,000.
Stipulate a pilot signal with sensitivity and specificity both $4/5$. Positive
and negative signals each occur with probability $1/2$. Conditional expected
construction values are USD~68,000 and minus USD~28,000, so the pilot policy
builds only after a positive signal. Expected value before pilot cost and delay
is USD~34,000. This also follows directly from
\[
 \tfrac12\tfrac45(100{,}000)
       +\tfrac12\tfrac15(-60{,}000)=34{,}000,
\]
which keeps the loss from an unfavorable state with a positive signal explicit.
At $K=10{,}000$ and $d=.9$, the pilot option is worth USD~20,600: only USD~600
more than immediate construction. Without delay it is worth USD~24,000.
Comparing the USD~34,000 gross value with a zero baseline would incorrectly
discard the already available USD~20,000 construction option.
\citep{exec14_l8_economics}

The example gives testable planning thresholds instead of a declared ROI.
At signal quality $4/5$, the pilot is weakly preferred to immediate construction exactly when
$d\geq15/17$. With $d=.9$ and symmetric signal quality $s$ in the regime
$s>5/8$, where the negative signal leads to declining construction, the gross
option value is $80{,}000s-30{,}000$. The break-even quality is
\[
 s\geq \frac{30{,}000+(20{,}000+K)/d}{80{,}000}
       =\frac{19}{24}\approx.7917.
\]
At lower quality the correct action after each signal must be recalculated;
the displayed branch cannot be extrapolated through a change of optimal action.
Perfect information would have gross value USD~50,000, an incremental
USD~30,000 over immediate construction before costs and delay. No imperfect
signal about this same fixed state and action set exceeds that perfect-information
value, because conditional maximization cannot outperform choosing with the
state observed. This comparison says nothing about pilots that themselves
create new feasible capabilities, which require a different model.
\citep{exec14_l8_economics}

For the causal proposal, the first signal should come from a bounded public or
generated prototype: can a selected adapter preserve the target; can it refuse
an unidentified request; does a proposed uncertainty method have useful width
under its declared law; and how much review and recovery does the complete
workflow require? These outcomes can change a construction decision without
claiming bank validation. A later authentic experiment becomes justified when
its expected information can resolve a consequential uncertainty that the
prototype cannot. The staged decision should include delay and the ability to
stop, as well as favorable technical results. \citep{exec14_l8_design}

\section{Transferring workflow value requires the deployment mixture}
\label{exec14:l8-transfer}

Even if a prototype measured paired workflow value exactly on its own cases,
the average need not transport to the requested bank population. Write
$\Delta_j$ for the incremental saving, after errors and full costs, in stratum
$j$, and $w_j$ for the deployment proportion. Under the strong assumption that
within-stratum values transport, deployment value is
\[
 \Delta_{\mathrm{deploy}}=\sum_j w_j\Delta_j.
\]
This assumption includes task difficulty, information availability, feasible
actions, review and losses within each stratum. Matching a topic name alone
does not establish it. If the weights are unknown but constrained to a set
$\mathcal W$, the corresponding value set is the range of the displayed linear
function over $\mathcal W$. Its extrema can be computed without inventing a
single deployment average. \citep{exec14_l8_economics}

Suppose the workflow saves USD~10 on a routine case and loses USD~30 on a
difficult case, with all components stipulated. If $w$ is the difficult-case
fraction, incremental value is $10-40w$ dollars per requested case. A public
development collection with $w=1/5$ has a USD~2 advantage, while a deployment
mix with $w=3/5$ has a USD~14 disadvantage. The threshold is $w=1/4$. Knowing
only that $w\in[.1,.4]$ gives a value range from minus USD~6 to plus USD~6.
Neither a public positive result nor uncertainty about the target mixture
justifies choosing its favorable endpoint. Selection by the assistant can
change the mix again; the routing policy and fallback outcomes must be part
of the transported object. \citep{exec14_l8_economics}

This is especially relevant when business translation is the proposed benefit.
A public dataset with clean treatment and outcome names can test numerical
orchestration yet supply little information about ambiguous requests, changing
population definitions or delayed records. A difficult-case stratum should be
defined by the causal task and information obstacle, not merely by a model's
low confidence. Language assistance can suggest the strata and identify missing
fields, but their prevalence and consequences need separate evidence. Feng
et al.'s reliability study restricts its accessible-corpus analysis and notes
its largely commonsense setting; it does not provide an error law for all
closed models or complex dynamic banking problems. That limitation supports
explicit transport assumptions, not an extrapolated bank failure rate.
\citep{v12_la_2407_19638}

Prices also need a transport argument. A prototype's observed analyst minutes
are resource use; converting them into savings requires a decision about
reassigned capacity, wage overhead and whether staffing expenditure changes.
Compute cost can have minimum charges, cached inputs, reserved capacity or
shared infrastructure. A model with lower variable cost can require more
adjudication or a larger fixed integration. For annual requested volume $n$,
eligible-use fraction $u$, per-used-request incremental value $v$ and additional
fixed cost $F$, the elementary no-discount scenario is $nuv-F$. Both $u$ and
$v$ can change with expansion, and $v$ already includes the fallback and loss
terms. A cost-per-answer claim should therefore be accompanied by the requested
denominator and resource totals. This proposal supplies the accounting and
break-even questions; actual deployment weights, prices and accepted values
are prospective measurements. \citep{exec14_l8_design}

\section{What each banking case still needs to justify further work}
\label{exec14:l8-research}

Funding has a concrete immediate prototype route: connect a selected summary
estimator to joint nuisance and dependence accounting, preserve the complete
response set, and propagate it through the stated profit map. Generated
calendars and known inclusion designs can reveal an invalid covariance or
misstated household unit cheaply. Public bank observations can establish
availability and preparation burden. Neither route by itself verifies a bank's
proxy exclusion, household truth or future margin. The economic question is
whether improved information or precision changes a treasury action enough to
justify the next acquisition. An interval that still spans consequential gains
and losses is an informative inconclusive result.

Servicing requires a policy that uses the information actually available at
each action, and an evaluation that preserves the original eligible cohort.
The finite-table example supplies an exact reference and a conservative bound
that is too wide for useful selection. A first prototype should compare a
specified estimator with that reference under controlled assignment, support
and dropout. Sharper valid uncertainty is a research objective; it is not
created by asking an agent to report a narrower interval. Hidden-state and
informative-observation examples with equal observed laws require an information
change or a defensible restriction. The eventual value of resolving them depends
on feasible servicing actions and rewards, including assistance costs.

Discovery requires a specified graph family, an honest identification state and
an effect or decision comparison after graph uncertainty. A prototype can first
check endpoint and node-order conversions against generated truth, then test a
selected discovery procedure under its own assumptions. The investment question
is whether a retained ambiguity changes the offer decision, and which additional
evidence can resolve it. Agreement among packages or among repeated language
answers is not an authenticated signal of the causal state. A stronger claim
about language-assisted discovery requires a separate controlled comparison of
the data, source evidence and expert effort supplied to each alternative.

Observation construction needs different evidence again. Valid clocks and
unchanged hashes can coexist with false labels, duplicated origins and wrong
inclusion probabilities. A public/generated prototype can test measurement
correction and the transmission of its uncertainty, including singular cases.
The agent's contribution may be to expose the missing audit or distinguish an
outcome clock from a feature clock. Authentic truth labels, ownership and joint
inclusion evidence are later inputs when the decision benefit warrants their
cost. The observation module should report which uncertainty remains instead
of awarding a truthful-population label from a successful timestamp check.
\citep{exec14_l8_design}

These routes have known engineering steps and genuine scientific limits. We
know how to implement a bounded adapter, preserve failed requests, compare exact
references, collect workflow costs and prepare a controlled pilot. We do not
know a universal theorem for arbitrary hidden-state decisions, unrestricted
dependent learners or arbitrary temporal graph effects, and no language model
can make observationally equivalent worlds distinguishable without additional
information. We also do not yet know the empirical accuracy, calibrated risk,
review burden or accepted bank value of the proposed assistance. A coherent
proposal treats each unknown as a specific next decision: prove a restricted
claim, change the information, test a prototype or decline the unsupported
use. Independent journeys can proceed under their own prerequisites. Bank
experiments follow demonstrated prototype merit and a useful information
question. \citep{exec14_l8_design}

The final chapter brings these conditions into one investment recommendation. The relevant comparison includes the option to retain an established analyst workflow, build a narrow reusable component, or defer an information-intensive capability. Each choice must be judged against the same decision and full cost, with the evidence that would reverse it stated explicitly.

\label{exec14:l8-development-end}
